现在提出一个数学问题, 其解将描述物体的平衡构形. 平衡由这类问题描述的材料称为\emph{弹性}材料. 首先定义\emph{能量泛函}:
\begin{equation}\label{eq:elasticity-total-energy-functional}
\mathcal{E}(\bm{p}):=\int_\Omega\mathcal{W}(\bm{x},\bm{F}(\bm{x}))\,\mathrm{d}\bm{x}+\int_\Omega\Psi(\bm{p}(\bm{x}))\,\mathrm{d}\bm{x}.
\end{equation}
这里 $\bm{F}:=\nabla\bm{p}$, 而 $\mathcal{W}:\Omega\times Q\to\mathbb{R}$ 为\emph{储能密度}, 其中 $Q$ 为具有正行列式的 $3\times3$ 矩阵的集合. 函数 $\Psi:\mathbb{R}^3\to\mathbb{R}$ 为\emph{势能密度}. 能量中的第一项称为\emph{储能泛函}, 它描述材料内部因机械形变而储存的能量. 第二项是\emph{势能泛函}, 表示来自外力 (假设其为保守力) 的能量.

对密度泛函作如下假设.\begin{enumerate}\item 假设 $\mathcal{W}\in C^2(\Omega\times Q)$ 且 $\Psi\in C^1(\mathbb{R}^3)$.\item 对某个 $p>3$, 存在常数 $k>0$ 和函数 $\omega\in C_b(\Omega)$, 使得
\begin{equation}\label{eq:elasticity-stored-energy-p-growth-bound}
\mathcal{W}(\bm{x},\bm{F})\geq\omega(\bm{x})+k|\bm{F}|^p.
\end{equation}
注意: 从现在起, $\mathcal{D}_{\mathcal{E}}$ 的定义中使用这一 $p$.\item 存在常数 $C$, 使对每个 $\bm{p}\in\mathcal{D}_{\mathcal{E}}$ 都有
\begin{equation}\label{eq:elasticity-potential-energy-uniform-lower-bound}
\int_\Omega\Psi(\bm{p}(\bm{x}))\,\mathrm{d}\bm{x}\geq C.
\end{equation}
\item 存在 $\bm{p}\in\mathcal{D}_{\mathcal{E}}$, 使 $\mathcal{E}(\bm{p})<\infty$.\item 对每个 $\bm{x}\in\Omega$, 当 $\det F\to0$ 时, 有 $\mathcal{W}(\bm{x},\bm{F})\to\infty$.\end{enumerate}

现在将物体的平衡构形定义为使能量 $\mathcal{E}$ 极小的构形; 即若
\begin{equation}\label{eq:elasticity-equilibrium-global-minimization}
\mathcal{E}(\overline{\bm{p}})\leq\mathcal{E}(\bm{p})\quad\text{for all }\bm{p}\in\mathcal{D}_{\mathcal{E}}.
\end{equation}
则称 $\overline{\bm{p}}\in\mathcal{D}_{\mathcal{E}}$ 为\emph{平衡态}或\emph{极小元}. 本节的首要目标是考察极小化问题 \eqref{eq:elasticity-equilibrium-global-minimization} 的解的存在性. 但在此之前, 我们希望揭示这一问题与椭圆偏微分方程的联系. 下面的结果将说明这一点.
